\documentclass[11pt,a4paper]{article}

% --- Paquetes ---
\usepackage[utf8]{inputenc}
\usepackage[T1]{fontenc}
\usepackage[spanish]{babel}
\usepackage{amsmath, amssymb, amsthm}
\usepackage{geometry}
\usepackage{tikz}
\usepackage{pgfplots}
\usepackage{enumitem}
\usepackage{hyperref}
\usepackage{xcolor}
\usepackage{tcolorbox}

\geometry{margin=2.5cm}
\pgfplotsset{compat=1.18}

% --- Definiciones ---
\newcommand{\barvar}[1]{\overline{#1}}

\newtcolorbox{definicionbox}[1][]{colback=yellow!10, colframe=red!70!black, title=#1, fonttitle=\bfseries}
\newtcolorbox{resultadobox}[1][]{colback=yellow!10, colframe=red!80!black, title=#1, fonttitle=\bfseries}
\newtcolorbox{explicacionbox}[1][]{colback=red!5, colframe=red!60!black, title=#1, fonttitle=\bfseries}

% --- Documento ---
\title{\textbf{Clase 2: Aspa Keynesiana} \\ \large Macroeconomía II}
\author{Universidad Católica del Norte}
\date{24 de agosto, 2026}

\begin{document}
\maketitle
\tableofcontents
\newpage

% ============================================================
\section{Notación Matemática}
% ============================================================

Durante el curso utilizaremos la siguiente tipología de variables.

\subsection{Tipos de Variables}

\begin{enumerate}[label=\alph*)]
    \item \textbf{Variables Endógenas:} Son determinadas \emph{dentro} del modelo.
    \begin{itemize}
        \item Ejemplo: En el Aspa Keynesiana serían $Y$ (ingreso) y $DA$ (demanda agregada).
        \item Notación: letra mayúscula. Ej: $Y$.
    \end{itemize}

    \item \textbf{Variables Exógenas:} Son determinadas \emph{fuera} del modelo.
    \begin{itemize}
        \item Ejemplo: Transferencias ($\barvar{Tr}$).
        \item Notación: letra mayúscula con una barra encima. Ej: $\barvar{Tr}$.
    \end{itemize}
\end{enumerate}

\subsection{Parámetros y Funciones}

\textbf{Parámetros:} Valores fijos que determinan la relación entre dos variables.
\begin{itemize}
    \item Notación: letra minúscula. Ej: $c$.
\end{itemize}

\textbf{Función:} Relación entre dos variables.
\begin{itemize}
    \item Notación: $I = I(r)$, donde el signo ``$=$'' se lee como ``igual'' (identidad funcional).
\end{itemize}

% ============================================================
\section{Aspa Keynesiana}
% ============================================================

\begin{itemize}
    \item Modelo simplificado.
    \item Nace para entender el \textbf{Mercado de Bienes y Servicios}.
    \item Supuestos:
    \begin{enumerate}
        \item Precios fijos.
        \item Economía cerrada.
        \item Corto plazo.
        \item Con Gobierno.
    \end{enumerate}
\end{itemize}

\subsection{Demanda Agregada}

\begin{definicionbox}[Demanda Agregada]
Valor total de todos los bienes y servicios que se \textbf{planean demandar} en la economía.
\end{definicionbox}

\subsubsection{Tipos de Bienes y Servicios que se demandan}

\paragraph{1) Consumo (C)}
Bienes y servicios demandados por los hogares.
\begin{itemize}
    \item \textbf{Bienes Duraderos:} Bienes que pueden ser utilizados durante un largo período. Ej: Auto.
    \item \textbf{Bienes No Duraderos:} Bienes que se consumen rápidamente. Ej: Comida.
    \item \textbf{Servicios:} No son bienes físicos pero son consumidos. Ej: Salud, Educación, Entretención, etc.
\end{itemize}

\paragraph{2) Inversión (I) --- Formación Bruta de Capital Fijo}
Demanda por bienes de capital. Los bienes de capital son insumos de producción que permiten generar mayor producción.
\begin{itemize}
    \item Ej: Una excavadora para una empresa constructora.
    \item Ej: Una impresora para una empresa editorial.
    \item Ej: Computadores para una empresa de servicios.
    \item \textbf{Construcciones:} Construcciones de nuevas fábricas, oficinas y tiendas.
    \item \textbf{Inventarios o Existencias:} Bienes que serán vendidos más adelante.
    \item \textbf{Inversión Pública:} Bienes de capital que demanda el Gobierno.
\end{itemize}

\paragraph{3) Gobierno (G)}
Bienes y servicios demandados por el Gobierno.
\begin{itemize}
    \item Salarios de los trabajadores del sector público.
    \item Administración.
    \item Servicios Públicos (Salud, Educación Pública, etc.).
\end{itemize}

% ============================================================
\subsection{Condición de Equilibrio del Mercado de Bienes y Servicios}
% ============================================================

\begin{resultadobox}[Condición de Equilibrio]
\begin{equation}
    DA = Y \tag{1}
\end{equation}
donde:
\begin{itemize}
    \item $DA$: Demanda Planeada (también llamada ``Gasto Planeado'' o ``Función de Demanda Agregada'').
    \item $Y$: Demanda Efectiva (también llamada ``Gasto Efectivo'').
\end{itemize}
\end{resultadobox}

Esta condición establece que el Mercado de Bienes y Servicios está en equilibrio cuando el volumen producido es igual a la cantidad demandada.

\subsubsection{Gasto Efectivo ($Y$)}

\begin{itemize}
    \item Gasto total efectivo de la economía.
    \item En Macroeconomía I aprendimos, por Contabilidad Nacional, que esto se puede medir de 3 formas:
    \begin{enumerate}
        \item \textbf{Método del Gasto Final:} $Y = C + I + G + XN$
        \item \textbf{Método del Valor Agregado.}
        \item \textbf{Método del Ingreso:} Suma de todos los ingresos (salarios, etc.) de la economía.
    \end{enumerate}
    \item La producción se mide a través del indicador PIB (Producto Interno Bruto).
    \item El gasto efectivo se mide por medio de la producción total de bienes y servicios de una economía en un año, usualmente a través del PIB.
\end{itemize}

\begin{tcolorbox}[colback=blue!5, colframe=blue!60!black]
Por lo tanto: $Y \Rightarrow$ Producción, Gasto Efectivo, Ingreso.
\end{tcolorbox}

\subsubsection{Gasto Planeado o Demanda Agregada ($DA$)}

Es lo que planean los hogares, las empresas y el Gobierno demandar a la economía, en términos de Bienes y Servicios, en un año.

$DA$ es una función, ya que lo que se planea demandar depende de varios factores.

\begin{resultadobox}[Esquema de Demanda Agregada]
\begin{equation}
    DA = C + I + G \tag{2}
\end{equation}
donde $C$ = Consumo, $I$ = Inversión, $G$ = Gasto Gubernamental.
\end{resultadobox}

Ahora empezamos a ver todas las funciones que componen el esquema de Demanda Agregada.

% ============================================================
\section{Funciones del Modelo}
% ============================================================

\subsection{Función de Consumo ($C$)}

Describe la relación entre consumo e ingreso:
\[
    C = f(Y)
\]
La relación entre consumo e ingreso es \textbf{positiva}. Escogemos la siguiente forma funcional:

\begin{resultadobox}[Función de Consumo]
\begin{equation}
    C = \barvar{C} + c \, Y_d \tag{3}
\end{equation}
con $\barvar{C} > 0$ y $0 < c < 1$.
\end{resultadobox}

Donde:
\begin{itemize}
    \item $\barvar{C}$ = \textbf{Consumo Autónomo}, es decir, todos los factores que afectan el consumo pero que no son el ingreso. Es positivo ($\barvar{C} > 0$).
    \item $c$ = \textbf{Propensión Marginal a Consumir}. Cambio en el consumo total por aumento de una unidad del Ingreso Disponible. $0 < c < 1$. Por ejemplo, $c = 0.8$ implica que la sociedad consume el 80\% del Ingreso Disponible.
\end{itemize}

\subsection{Ingreso Disponible ($Y_d$)}

El ingreso disponible es igual al ingreso que le queda a los agentes de la economía una vez que le restan los impuestos y se le suma cualquier tipo de transferencia.

\begin{resultadobox}[Ingreso Disponible]
\begin{equation}
    Y_d = Y - \barvar{T} + \barvar{Tr} \tag{4}
\end{equation}
\end{resultadobox}

Donde:
\begin{itemize}
    \item $\barvar{T}$ = Impuesto suma fija.
    \item $\barvar{Tr}$ = Transferencias (Ej: Bonos, etc.). Las transferencias son sumas de dinero otorgadas a las familias sin un trabajo productivo a cambio.
\end{itemize}

\subsubsection{Consumo en función de $Y$}

Reemplazando (4) en (3):
\begin{align}
    C &= \barvar{C} + c\left(Y - \barvar{T} + \barvar{Tr}\right) \notag \\
    C &= \barvar{C} + cY - c\barvar{T} + c\barvar{Tr} \notag
\end{align}

\begin{resultadobox}[Función de Consumo en función de $Y$]
\begin{equation}
    C = \underbrace{\barvar{C}}_{\text{Endógena}} - \underbrace{c\barvar{T} + c\barvar{Tr}}_{\text{Exógenas}} + \underbrace{cY}_{\text{Endógena}} \tag{5}
\end{equation}
\end{resultadobox}

\subsubsection{¿Cómo graficar la función de consumo?}

\begin{enumerate}
    \item \textbf{1\textsuperscript{er} Paso:} Identificar variables endógenas vs exógenas. Las variables endógenas serán los ejes del gráfico: $C$, $Y$.
    \item \textbf{2\textsuperscript{do} Paso:} Identificar si la pendiente es positiva o negativa. La pendiente de la función de consumo es $c > 0$, positiva. Entre mayor ingreso existe, la gente más consume.
    \item \textbf{3\textsuperscript{er} Paso:} Identificar los puntos de corte en los ejes.
\end{enumerate}

La función es $C = \barvar{C} - c\barvar{T} + c\barvar{Tr} + cY$.

\begin{itemize}
    \item Si $Y = 0$: $C = \barvar{C} - c\barvar{T} + c\barvar{Tr} > 0$ (positivo).
    \item Si $C = 0$: $Y = -\dfrac{\barvar{C} - c\barvar{T} + c\barvar{Tr}}{c}$ (negativo).
\end{itemize}

\begin{center}
\begin{tikzpicture}
    \begin{axis}[
        axis lines=middle,
        xlabel={$Y$}, ylabel={$C$},
        xmin=-3, xmax=8, ymin=-1, ymax=6,
        xtick=\empty, ytick=\empty,
        every axis x label/.style={at={(ticklabel* cs:1)}, anchor=west},
        every axis y label/.style={at={(ticklabel* cs:1)}, anchor=south},
        width=10cm, height=7cm,
    ]
    % Función de consumo
    \addplot[domain=-2:7, thick, violet] {1.5 + 0.6*x};
    \node[right, violet] at (axis cs:5.5,5) {$C = \barvar{C} - c\barvar{T} + c\barvar{Tr} + cY$};

    % Interceptos
    \node[left, teal] at (axis cs:0,1.5) {$\barvar{C} - c\barvar{T} + c\barvar{Tr}$};
    \node[below] at (axis cs:-2.5,0) {$-\frac{\barvar{C} - c\barvar{T} + c\barvar{Tr}}{c}$};

    % Pendiente
    \node[below right, violet] at (axis cs:2,2.7) {Pendiente $= c$};

    \end{axis}
\end{tikzpicture}
\end{center}

\subsection{Consumo y Ahorro}

¿Qué ocurre con $(1-c)$? Es la parte del Ingreso Disponible que no se consume. Si no se gasta, se ahorra.

\begin{resultadobox}[Restricción Presupuestaria]
\begin{equation}
    S = Y - C \tag{6}
\end{equation}
donde $S$ = Ahorro.
\end{resultadobox}

Sustituyendo (5) en (6):
\begin{align}
    S &= Y - \left(\barvar{C} - c\barvar{T} + c\barvar{Tr} + cY\right) \notag \\
    S &= Y - \barvar{C} + c\barvar{T} - c\barvar{Tr} - cY \notag \\
    S &= -\barvar{C} + c\barvar{T} - c\barvar{Tr} + \underbrace{(1-c)}_{s}Y \notag
\end{align}

donde $s = (1-c)$ es la \textbf{propensión marginal a ahorrar}.

\subsection{Inversión ($I$)}

Por el momento asumiremos que la inversión es exógena:

\begin{resultadobox}[Inversión]
\begin{equation}
    I = \barvar{I} \tag{7}
\end{equation}
\end{resultadobox}

\subsection{Gasto Gubernamental ($G$)}

Se asume que el gasto de gobierno es exógeno:

\begin{resultadobox}[Gasto Gubernamental]
\begin{equation}
    G = \barvar{G} \tag{8}
\end{equation}
\end{resultadobox}

% ============================================================
\section{Función de Demanda Agregada}
% ============================================================

Para calcular la función de Demanda Agregada, sustituimos (5), (7), y (8) en (2):
\begin{align}
    DA &= C + I + G \notag \\
    DA &= \underbrace{\barvar{C} - c\barvar{T} + c\barvar{Tr} + \barvar{I} + \barvar{G}}_{\text{Variables Exógenas}} + \underbrace{cY}_{\text{Variable Endógena}} \notag
\end{align}

Todas las variables exógenas se pueden agrupar en un solo término que denominamos \textbf{Gasto Autónomo}:

\begin{resultadobox}[Función de Demanda Agregada]
\begin{equation}
    DA = \barvar{A}_0 + cY \tag{9}
\end{equation}
donde:
\[
    \barvar{A}_0 = \barvar{C} - c\barvar{T} + c\barvar{Tr} + \barvar{I} + \barvar{G}
\]
\end{resultadobox}

La \textbf{Función de Demanda Agregada} describe la relación entre el ingreso total de una economía y la demanda agregada.

% ============================================================
\section{Determinación de la Producción de Equilibrio}
% ============================================================

El nivel de producción en el cual una economía está en equilibrio se encuentra cuando el Gasto Efectivo es igual al Gasto Planeado:
\begin{align}
    DA &= Y \quad \rightarrow \text{Ecuación (1)} \notag \\
    DA &= \barvar{A}_0 + cY \quad \rightarrow \text{Ecuación (9)} \notag
\end{align}

Reemplazamos (9) en (1):
\begin{align}
    \barvar{A}_0 + cY &= Y \notag \\
    Y - cY &= \barvar{A}_0 \notag \\
    (1-c)Y &= \barvar{A}_0 \notag
\end{align}

\begin{resultadobox}[Nivel de Producción de Equilibrio]
\begin{equation}
    Y^* = \frac{\barvar{A}_0}{1 - c}
\end{equation}
\end{resultadobox}

\subsection{Gráfico del Aspa Keynesiana}

Para graficar:
\begin{enumerate}
    \item \textbf{Variables endógenas:} $DA = Y \Rightarrow DA, Y$; \quad $DA = \barvar{A}_0 + cY \Rightarrow DA, Y$.
    \item \textbf{Pendiente:} $c > 0$ (positiva).
    \item \textbf{Puntos de corte:}
    \begin{itemize}
        \item Si $Y=0$: $DA = \barvar{A}_0$.
        \item Si $DA = 0$: $Y = -\barvar{A}_0 / c$.
    \end{itemize}
\end{enumerate}

\begin{center}
\begin{tikzpicture}
    \begin{axis}[
        axis lines=middle,
        xlabel={$Y$}, ylabel={$DA$},
        xmin=-1, xmax=10, ymin=-1, ymax=10,
        xtick=\empty, ytick=\empty,
        every axis x label/.style={at={(ticklabel* cs:1)}, anchor=west},
        every axis y label/.style={at={(ticklabel* cs:1)}, anchor=south},
        width=11cm, height=9cm,
    ]
    % Línea de 45 grados (DA = Y)
    \addplot[domain=0:9, thick, blue] {x};
    \node[above left, blue] at (axis cs:8.5,8.5) {$DA = Y$};

    % Función DA = Ao + cY (con c = 0.6 y Ao = 2)
    \addplot[domain=0:9, thick, red!70!black] {2 + 0.6*x};
    \node[right, red!70!black] at (axis cs:7.5,6.5) {$DA = \barvar{A}_0 + cY$};

    % Punto de equilibrio E
    \addplot[only marks, mark=*, mark size=3pt, red!70!black] coordinates {(5, 5)};
    \node[above left] at (axis cs:5,5) {$E$};

    % Líneas punteadas al equilibrio
    \addplot[dashed, gray] coordinates {(5,0) (5,5)};
    \addplot[dashed, gray] coordinates {(0,5) (5,5)};

    % Etiquetas
    \node[left] at (axis cs:0,2) {$\barvar{A}_0$};
    \node[left] at (axis cs:0,5) {$DA^* = \dfrac{\barvar{A}_0}{1-c}$};
    \node[below] at (axis cs:5,0) {$Y^* = \dfrac{\barvar{A}_0}{1-c}$};

    % Pendiente label
    \node[below right, red!70!black] at (axis cs:6.5,5.8) {Pendiente $= c$};

    % Ángulo de 45 grados
    \draw[blue, thick] (axis cs:1.2,0) arc[start angle=0, end angle=45, radius=1cm];
    \node[above right, blue] at (axis cs:1.3,0.3) {$45°$};

    \end{axis}
\end{tikzpicture}
\end{center}

% ============================================================
\subsection{Estática Comparativa: ¿Qué ocurre si aumenta $c$?}
% ============================================================

Si la propensión marginal a consumir aumenta (por ejemplo, de $c_1 = 0.6$ a $c_2 = 0.8$), la sociedad se vuelve más consumista. ¿Qué ocurre con la producción de equilibrio?

\[
    Y_1^* = \frac{\barvar{A}_0}{1-c_1} \quad \longrightarrow \quad Y_2^* = \frac{\barvar{A}_0}{1-c_2}
\]

Como $c_2 > c_1$, entonces $(1-c_2) < (1-c_1)$, por lo que $Y_2^* > Y_1^*$.

\begin{center}
\begin{tikzpicture}
    \begin{axis}[
        axis lines=middle,
        xlabel={$Y$}, ylabel={$DA$},
        xmin=-1, xmax=12, ymin=-1, ymax=12,
        xtick=\empty, ytick=\empty,
        every axis x label/.style={at={(ticklabel* cs:1)}, anchor=west},
        every axis y label/.style={at={(ticklabel* cs:1)}, anchor=south},
        width=11cm, height=9cm,
    ]
    % Línea de 45 grados
    \addplot[domain=0:11, thick, blue] {x};
    \node[above left, blue] at (axis cs:10,10) {$DA = Y$};

    % DA con c1
    \addplot[domain=0:11, thick, orange] {2 + 0.6*x};
    \node[right, orange] at (axis cs:9,7.4) {$DA = \barvar{A}_0 + c_1 Y$};

    % DA con c2 (pivote desde Ao)
    \addplot[domain=0:11, thick, red] {2 + 0.8*x};
    \node[right, red] at (axis cs:8.5,8.8) {$DA = \barvar{A}_0 + c_2 Y$};

    % Equilibrio E1
    \addplot[only marks, mark=*, mark size=3pt, violet] coordinates {(5, 5)};
    \node[below left, violet] at (axis cs:5,5) {$E_1$};

    % Equilibrio E2
    \addplot[only marks, mark=*, mark size=3pt, red] coordinates {(10, 10)};
    \node[above left, red] at (axis cs:10,10) {$E_2$};

    % Líneas punteadas
    \addplot[dashed, violet] coordinates {(5,0) (5,5)};
    \addplot[dashed, red] coordinates {(10,0) (10,10)};

    % Etiquetas en eje X
    \node[below, violet] at (axis cs:5,0) {$Y_1^*$};
    \node[below, red] at (axis cs:10,0) {$Y_2^*$};

    % Flecha de pivote
    \draw[->, thick, gray, dashed] (axis cs:7,6.2) to[bend left=20] (axis cs:7,7.6);
    \node[right, gray] at (axis cs:7.2,6.9) {Pivote};

    \end{axis}
\end{tikzpicture}
\end{center}

\begin{explicacionbox}[Explicación]
Si la sociedad comienza a consumir más, la demanda agregada aumenta ($\uparrow DA$). La producción responde disminuyendo inventarios y contratando a trabajadores para alcanzar la nueva demanda ($\uparrow Y$). Esto genera un nuevo aumento en el ingreso y por ende en el consumo ($\uparrow C$) y la $DA$ ($\uparrow DA$), hasta que eventualmente se llega a un nuevo equilibrio.
\end{explicacionbox}

% ============================================================
\subsection{Estática Comparativa: ¿Qué ocurre si aumenta $\barvar{A}_0$?}
% ============================================================

Pensar en las razones por las cuales $\barvar{A}_0$ podría aumentar (aumento en $\barvar{G}$, $\barvar{I}$, $\barvar{C}$, $\barvar{Tr}$, o disminución de $\barvar{T}$).

Si $\barvar{A}_0^1 < \barvar{A}_0^2$:

\[
    Y_1^* = \frac{\barvar{A}_0^1}{1-c} \quad \longrightarrow \quad Y_2^* = \frac{\barvar{A}_0^2}{1-c}
\]

\begin{center}
\begin{tikzpicture}
    \begin{axis}[
        axis lines=middle,
        xlabel={$Y$}, ylabel={$DA$},
        xmin=-1, xmax=12, ymin=-1, ymax=12,
        xtick=\empty, ytick=\empty,
        every axis x label/.style={at={(ticklabel* cs:1)}, anchor=west},
        every axis y label/.style={at={(ticklabel* cs:1)}, anchor=south},
        width=11cm, height=9cm,
    ]
    % Línea de 45 grados
    \addplot[domain=0:11, thick, blue] {x};
    \node[above left, blue] at (axis cs:10,10) {$DA = Y$};

    % DA con Ao1
    \addplot[domain=0:11, thick, orange] {2 + 0.6*x};
    \node[right, orange] at (axis cs:9,7.4) {$DA = \barvar{A}_0^1 + cY$};

    % DA con Ao2 (desplazamiento paralelo)
    \addplot[domain=0:11, thick, red] {4 + 0.6*x};
    \node[right, red] at (axis cs:8,8.8) {$DA = \barvar{A}_0^2 + cY$};

    % Equilibrio E1
    \addplot[only marks, mark=*, mark size=3pt, violet] coordinates {(5, 5)};
    \node[below left, violet] at (axis cs:5,5) {$E_1$};

    % Equilibrio E2
    \addplot[only marks, mark=*, mark size=3pt, orange!80!red] coordinates {(10, 10)};
    \node[above left, orange!80!red] at (axis cs:10,10) {$E_2$};

    % Líneas punteadas
    \addplot[dashed, violet] coordinates {(5,0) (5,5)};
    \addplot[dashed, orange!80!red] coordinates {(10,0) (10,10)};

    % Etiquetas
    \node[left] at (axis cs:0,2) {$\barvar{A}_0^1$};
    \node[left] at (axis cs:0,4) {$\barvar{A}_0^2$};
    \node[below, violet] at (axis cs:5,0) {$Y_1^*$};
    \node[below, orange!80!red] at (axis cs:10,0) {$Y_2^*$};

    % Flecha de desplazamiento
    \draw[->, thick, orange!70!red, dashed] (axis cs:3,3.8) -- (axis cs:3,5.8);
    \node[right, orange!70!red] at (axis cs:3.1,4.8) {Desplazamiento};

    % Recuadro
    \node[draw=red, thick, fill=red!5, anchor=west] at (axis cs:8,4) {$\barvar{A}_0^1 < \barvar{A}_0^2$};

    \end{axis}
\end{tikzpicture}
\end{center}

Un aumento en el gasto autónomo produce un \textbf{desplazamiento paralelo} hacia arriba de la función de demanda agregada, generando un nuevo equilibrio con mayor producción.

\end{document}
